\documentclass[11pt]{article}
\usepackage[margin=0.9in]{geometry}
\usepackage{amsmath,amssymb,amsthm}
\usepackage{graphicx}
\usepackage{hyperref}
\usepackage{microtype}
\hypersetup{hidelinks}

\newtheorem{theorem}{Theorem}
\newtheorem{lemma}[theorem]{Lemma}
\newcommand{\R}{\mathbb R}
\newcommand{\cH}{\mathcal H}
\newcommand{\sd}{\sigma_{\rm d}}
\newcommand{\one}{\mathbf 1}

\title{Strict Positivity and Simplicity of the Ground State\\
for a Functional-Difference Operator}
\author{Candidate full solution packet for arXiv:2109.05465}
\date{11 August 2026}

\begin{document}
\maketitle

\begin{center}
\fbox{\parbox{0.91\textwidth}{\textbf{Status.}
Candidate full affirmative answer; likely valid, pending human review.  Under
the hypothesis in the source problem, the lowest eigenvalue is simple and its
eigenfunction has a continuous representative that is strictly positive at
every real point.}}
\end{center}

\section{Source question}

Let $P=-i\,d/dx$, $b>0$, and
\[
 W_0(b)=2\cosh(bP),\qquad W_V(b)=W_0(b)-V
\]
on $L^2(\R)$, where $0\le V\in L^1(\R)$ and $W_V(b)$ is defined by its
Friedrichs form.  Laptev and Schimmer ask after Remark~4.1 of
\emph{A Sharp Lieb--Thirring Inequality for Functional Difference Operators},
arXiv:2109.05465 [1]:

\begin{quote}
Assume that $\sd(W_V(b))\subset[-2,2)$.  Is it true that the eigenfunction
corresponding to the lowest eigenvalue could be chosen strictly positive?
\end{quote}

\begin{figure}[h]
\centering
\includegraphics[width=\textwidth]{figures/open_problem_crop.png}
\caption{The exact question and its motivating rank-one calculation on source
PDF page~9.}
\end{figure}

We give an affirmative answer and also prove simplicity.  The case $V=0$ has
no discrete eigenvalue and is vacuous, so below $V$ is assumed nonzero.

\section{Proof intuition}

The obstruction seen by the source below energy $-2$ is precisely the loss of
positivity of the free Green function.  On the entire interval $[-2,2)$,
including the endpoint $-2$, that Green function is instead strictly positive.
Consequently the Birman--Schwinger operator at the lowest eigenvalue is a
compact positivity-improving integral operator on the support of $V$.

The only remaining point is to know which of its eigenvalues equals one.
Birman--Schwinger counting says that an eigenvalue above one would produce an
eigenvalue of $W_V(b)$ strictly below the alleged ground state.  Thus one is
the spectral radius.  A short Perron argument makes its eigenspace
one-dimensional with a strictly positive eigenvector.  Applying the free
resolvent to that vector reconstructs an $L^2$ eigenfunction which is a
convolution of two nonnegative nonzero functions, one of them everywhere
strictly positive.

\section{Statement}

\begin{theorem}\label{thm:main}
Let $b>0$ and let $0\le V\in L^1(\R)$ be nonzero.  Suppose
\[
 \sd(W_V(b))\subset[-2,2).
\]
Then $W_V(b)$ has a lowest discrete eigenvalue $\lambda_0$.  This eigenvalue
is simple, and its eigenspace is spanned by a function having a continuous
representative $\psi_0$ such that
\[
 \psi_0(x)>0\qquad(x\in\R).
\]
In particular, the source question has an affirmative answer.
\end{theorem}

\section{The positive free resolvent}

The Fourier symbol of $W_0(b)$ is $2\cosh(2\pi bk)$.  For $\lambda<2$ write
$R_\lambda=(W_0(b)-\lambda)^{-1}$.

\begin{lemma}\label{lem:green}
If $-2\le\lambda<2$, then $R_\lambda$ is convolution with a real, even
function $G_\lambda\in C_0(\R)\cap L^2(\R)$ which is strictly positive at
every point.
\end{lemma}

\begin{proof}
For $-2<\lambda<2$, write $\lambda=-2\cos\omega$ with $0<\omega<\pi$.
The source computes
\begin{equation}\label{eq:G}
 G_\lambda(t)=\frac{1}{2b\sin\omega}
 \frac{\sinh(\omega t/b)}{\sinh(\pi t/b)},
 \qquad
 G_\lambda(0)=\frac{1}{2\pi b}\frac{\omega}{\sin\omega}.
\end{equation}
For $t\ne0$, the two hyperbolic sines have the same sign.  Thus
$G_\lambda(t)>0$, and the displayed diagonal value is positive.  The quotient
is continuous at zero and decays exponentially at infinity because
$\omega<\pi$, proving the asserted regularity and integrability.

At $\lambda=-2$, passage to the limit $\omega\downarrow0$ gives
\begin{equation}\label{eq:Gendpoint}
 G_{-2}(t)=\frac{t}{2b^2\sinh(\pi t/b)},
 \qquad G_{-2}(0)=\frac{1}{2\pi b}.
\end{equation}
This function has the same properties and is strictly positive everywhere.
\end{proof}

Put $S=\{x:V(x)>0\}$, modulo null sets, and use $L^2(S)$ as the natural
nonzero support space.  Define
\[
 K_\lambda=V^{1/2}R_\lambda V^{1/2}.
\]

\begin{lemma}\label{lem:K}
For $-2\le\lambda<2$, $K_\lambda$ is a nonzero compact self-adjoint
nonnegative operator on $L^2(S)$.  It is positivity improving: if
$0\ne f\in L^2(S)$ and $f\ge0$, then $K_\lambda f>0$ almost everywhere on
$S$.
\end{lemma}

\begin{proof}
Its integral kernel is
\[
 k_\lambda(x,y)=V(x)^{1/2}G_\lambda(x-y)V(y)^{1/2}.
\]
Lemma~\ref{lem:green} makes $G_\lambda$ bounded, so
\[
 \|K_\lambda\|_{\rm HS}^2
 \le \|G_\lambda\|_\infty^2\|V\|_1^2<\infty.
\]
Thus $K_\lambda$ is compact.  It is self-adjoint because $G_\lambda$ is real
and even, and it is nonnegative in the Hilbert-space sense because
$R_\lambda\ge0$.  Finally, $k_\lambda(x,y)>0$ for almost every $(x,y)\in
S\times S$.  Integrating it against a nonzero nonnegative $f$ proves
positivity improvement and also nontriviality.
\end{proof}

We record the elementary Perron fact needed below, including its proof so that
no irreducibility theorem is hidden in the argument.

\begin{lemma}\label{lem:perron}
Let $K$ be a nonzero compact self-adjoint nonnegative integral operator on a
real or complex $L^2$-space.  If its kernel is strictly positive almost
everywhere, then $r=\|K\|$ is a simple eigenvalue and it has an eigenvector
$\phi>0$ almost everywhere.
\end{lemma}

\begin{proof}
Compactness and nonnegativity show that $r>0$ is an eigenvalue.  Since the
kernel is real, choose a nonzero real eigenvector $u$.  Pointwise kernel
positivity gives
\[
 \langle |u|,K|u|\rangle\ge \langle u,Ku\rangle=r\|u\|_2^2.
\]
The reverse inequality follows from $K\le rI$.  Hence equality holds and
$|u|$ is again an $r$-eigenvector.  Positivity improvement gives
$K|u|>0$ almost everywhere, so $|u|>0$ almost everywhere.

If the real $r$-eigenspace had dimension at least two, it would contain a
nonzero $w$ orthogonal to the positive eigenvector.  Both $w_+$ and $w_-$
would then be nonzero.  Strict positivity of the kernel gives
\[
 \langle |w|,K|w|\rangle-\langle w,Kw\rangle
 =4\langle w_+,Kw_-\rangle>0,
\]
contradicting $\langle |w|,K|w|\rangle\le r\|w\|_2^2=
\langle w,Kw\rangle$.  Thus the real eigenspace is one-dimensional, and the
same follows for its complexification.
\end{proof}

\section{Birman--Schwinger counting at the ground state}

We use the standard Birman--Schwinger principle in the same form setting as
the source paper.  For every $\lambda<2$,
\begin{equation}\label{eq:BS-count}
 \#\{\text{eigenvalues of }W_V(b)\text{ below }\lambda\}
 =\#\{\text{eigenvalues of }K_\lambda\text{ above }1\},
\end{equation}
with multiplicity.  Moreover,
\begin{equation}\label{eq:BS-eq}
 \ker(W_V(b)-\lambda)\longleftrightarrow\ker(K_\lambda-I)
\end{equation}
is a linear isomorphism, initially given from left to right by
$\psi\mapsto V^{1/2}\psi$.

For completeness, these statements follow by writing the form of
$W_V(b)-\lambda$ through $(W_0(b)-\lambda)^{1/2}$ and applying the min--max
principle to $I-R_\lambda^{1/2}VR_\lambda^{1/2}$.  The latter operator and
$K_\lambda$ have the same nonzero eigenvalues.  At eigenvalue one, the inverse
map in \eqref{eq:BS-eq} is
\begin{equation}\label{eq:reconstruct}
 \phi\longmapsto R_\lambda V^{1/2}\phi.
\end{equation}
This is also the Birman--Schwinger principle invoked in Section~3.1 of [1].

\begin{proof}[Proof of Theorem~\ref{thm:main}]
Because $V\ge0$ is nonzero, Theorem~1.3 of [1] gives at least one
eigenvalue below $2$.  The source also establishes that the essential spectrum
is $[2,\infty)$ and that the spectrum below $2$ consists of isolated
finite-multiplicity eigenvalues.  Such eigenvalues can accumulate only at
$2$, so the stated inclusion in $[-2,2)$ gives a lowest one, denoted
$\lambda_0$.

The correspondence \eqref{eq:BS-eq} says that $1$ is an eigenvalue of
$K_{\lambda_0}$.  Since no eigenvalue of $W_V(b)$ lies below $\lambda_0$,
\eqref{eq:BS-count} says that $K_{\lambda_0}$ has no eigenvalue strictly above
$1$.  Lemma~\ref{lem:K} makes this a nonnegative compact operator, and hence
\[
 \|K_{\lambda_0}\|=1.
\]
Lemmas~\ref{lem:K} and \ref{lem:perron} provide a unique (up to scalar)
eigenvector $\phi\in L^2(S)$ for eigenvalue one, normalized so that
$\phi>0$ almost everywhere on $S$.

Set $h=V^{1/2}\phi$, extended by zero off $S$.  Since
$V^{1/2},\phi\in L^2(S)$, Cauchy--Schwarz gives $h\in L^1(\R)$, and $h$ is
nonnegative and nonzero.  Define
\[
 \psi=R_{\lambda_0}h=G_{\lambda_0}*h.
\]
Lemma~\ref{lem:green} and Young's inequality give $\psi\in L^2(\R)$.  The
Birman--Schwinger equation yields
\[
 V^{1/2}\psi=K_{\lambda_0}\phi=\phi,
 \qquad V\psi=V^{1/2}\phi=h.
\]
Thus \eqref{eq:reconstruct}, or equivalently the quadratic-form identity,
shows that $W_V(b)\psi=\lambda_0\psi$.

There is no domain gap in this reconstruction.  Indeed $\widehat h$ is
bounded and
\[
 \int_\R
 \frac{2\cosh(2\pi bk)}{(2\cosh(2\pi bk)-\lambda_0)^2}
 |\widehat h(k)|^2\,dk<\infty,
\]
so $\psi$ belongs to the form domain of $W_0(b)$.  The equality
$V\psi=h\in L^1$ then gives the eigenvalue identity first in the form sense;
the representation theorem places $\psi$ in the operator domain.

Finally, $G_{\lambda_0}\in C_0(\R)$ and $h\in L^1(\R)$, so the convolution
is continuous.  For every $x\in\R$,
\[
 \psi(x)=\int_\R G_{\lambda_0}(x-y)h(y)\,dy>0,
\]
because the Green function is everywhere strictly positive and $h$ is
nonnegative and nonzero.  Simplicity of the eigenvalue one in
Lemma~\ref{lem:perron}, together with the eigenspace isomorphism
\eqref{eq:BS-eq}, proves simplicity of $\lambda_0$.
\end{proof}

\section{Verification, scope, and novelty}

The proof covers the endpoint $\lambda_0=-2$ explicitly through
\eqref{eq:Gendpoint}; no limiting eigenfunction argument is used.  It uses
only $V\in L^1$, not boundedness or strict positivity of $V$.  Positivity is
proved on all of $\R$, not merely almost everywhere, for the continuous
convolution representative.

Bounded local-index, parsed-arXiv, exact-phrase, title/author, DOI/citation,
and author-page searches were run on 11 August 2026.  They found the source
and adjacent 2025 nonselfadjoint work, but no later paper stating this answer.
Novelty confidence is moderate; this is not a priority claim.

\textbf{Human-review recommendation.}  Check form Birman--Schwinger counting
for $V\in L^1$, the endpoint kernel \eqref{eq:Gendpoint}, and reconstruction.
The Perron argument is self-contained; no unproved dependency remains.

\section*{Reference}
\small
\noindent [1] A.~Laptev and L.~Schimmer,
\emph{A Sharp Lieb--Thirring Inequality for Functional Difference Operators},
SIGMA \textbf{17} (2021), 105, 10 pp.; arXiv:2109.05465;
\href{https://doi.org/10.3842/SIGMA.2021.105}{doi:10.3842/SIGMA.2021.105}.

\end{document}
